% Hexweave: bounded-residency continual training of a paged byte-level language model on a CPU.
% Build: see README.md in this directory.
\documentclass[10pt]{article}

\usepackage[T1]{fontenc}
\usepackage[utf8]{inputenc}
\usepackage{lmodern}
\usepackage[margin=1in]{geometry}
\usepackage{amsmath,amssymb}
\usepackage{microtype}
\usepackage{booktabs}
\usepackage{xcolor}
\usepackage{caption}
\usepackage{enumitem}
\usepackage{multicol}
\usepackage[numbers,sort&compress]{natbib}
\usepackage{tikz}
\usetikzlibrary{shapes.geometric,positioning,arrows.meta,fit,calc}
\usepackage{pgfplots}
\pgfplotsset{compat=1.17}
\usepackage[colorlinks=true,linkcolor=black!70!blue,citecolor=black!60!green,urlcolor=black!70!blue]{hyperref}
\hypersetup{pdftitle={Training a Paged Language Model under a Fixed Memory Budget},pdfauthor={Jade Zaher}}

% Values that depend on the long run still in progress (Section 5.4). Fill these in when it ends.
\newcommand{\pending}[1]{\textcolor{red!75!black}{\fbox{\textsf{\footnotesize TODO: #1}}}}

\newcommand{\bpb}{bits/byte}
\captionsetup{font=small}
\setlength{\textfloatsep}{12pt plus 2pt minus 4pt}
\setlength{\floatsep}{10pt plus 2pt minus 2pt}
\newcommand{\nohyphens}{\hyphenpenalty=10000\exhyphenpenalty=10000\relax}
\setlength{\bibsep}{1pt plus 1pt}
\renewcommand{\bibfont}{\footnotesize}
\newcommand{\code}[1]{\texttt{#1}}

\definecolor{decaycol}{RGB}{200,90,30}
\definecolor{restartcol}{RGB}{30,90,180}

\title{Training a Paged Language Model under a Fixed Memory Budget:\\
Bounded-Residency, Per-Tile Continual Updates on a Desktop CPU}
\author{Jade Zaher}
\date{November 2026 \\ {\small Working paper; one small corpus; a second seed for one recipe only.}}

\begin{document}
\maketitle

\begin{abstract}
We want to train a language model whose parameter count is not capped by RAM: the user chooses a
memory budget, the model's parameters and optimizer state live on disk in independent blocks, and
learning is online and local, so that each training position updates only the blocks it used. We
accept slower training and some loss of per-step quality in exchange for bounded residency and
real-time updating. This paper reports the first measurements of such a system, Hexweave's
\emph{connected} path, on a commodity desktop CPU. A byte-level model of 2,051,840 coefficients is
split into seven \emph{tiles} on a hexagonal grid; a learned router sends each prediction through at
most four tiles; tiles and their Adam state are paged between an SSD and a bounded resident set; and
each training position takes one exact-gradient Adam step on the tiles its route visited. On
TinyStories (seed 7 unless stated, 1,024 held-out bytes per evaluation), a 24-hour run with a global
inverse-square-root rate decay plateaued between the byte bigram and unigram levels (windowed means
3.35--3.41~\bpb) while writing 3.70\,TB to disk. Replacing eager eviction with a write-behind buffer
and visit-frequency eviction left the learning curve unchanged and cut bytes written by about 75$\times$
in 30-minute pilots. Replacing the global decay with per-tile warm restarts triggered by accumulated
parameter change produced a run whose windowed held-out mean fell to 2.65~\bpb{} at 1.0--1.2M training
positions, below the matched-exposure bigram (3.31), with about 50\,MB peak working set; it ended at
its 24-hour budget after 1,691,199 positions with a last evaluation of 2.563 (best 2.476) and about
11\,GB written. A 6-hour run with a constant tile rate, no restarts and no forgetting matched or
slightly beat it in every completed 100k-position window up to 500k, so at this scale and horizon the
gain came from not letting the tile rate decay, not from restarts. A continuous plasticity law based
on the Lambert $W$ function led the constant rate by 0.06--0.125~\bpb{} at 200--400k positions but
also ran at a higher rate, so that lead is not attributed. A second seed of the restart recipe
followed the same curve with a slower start. In short: bounded-residency, per-tile online training
works at this small scale; the key ingredient was keeping the tile rate up; paging cost depends on
the store policy.
\end{abstract}

\section{Introduction}

\paragraph{Goal.}
The parameter count of a model trained on a single machine is usually bounded by memory: weights,
gradients and optimizer moments must all be addressable during each step. Offloading systems move
some of that state to host memory or NVMe~\citep{ren2021zerooffload,rajbhandari2021zeroinfinity}, but
they still perform a global step over every parameter for every batch, so traffic and time scale
with the whole model. We ask a different question: \emph{can a language model be trained on
commodity hardware under a memory budget the user chooses, independent of its total parameter
count, with updates that are local (touching only the parameters a training position actually used)
and continual (one position at a time, no epochs or global synchronization)?} We are willing to pay
for this in training time and in some per-step quality. In exchange we want two properties that
conventional training does not offer: resident memory that stays bounded as the model grows, and a
model that can keep learning in real time from the text it sees.

\paragraph{Approach.}
Hexweave splits a byte-level model into \emph{tiles}, small self-contained blocks with their own
attention, feed-forward network, router and optimizer state, laid out on an axial hexagonal grid. A
prediction starts at the centre tile and a learned router chooses, after each visit, a neighbouring
tile or a halt, for at most four visits. Because a prediction touches at most four tiles, only those
tiles need to be in memory for it, and only those tiles are updated by it. The \emph{connected} path
(the subject of this paper) keeps at most $R$ decoded tiles in RAM and pages the rest, with their
Adam moments, to and from an SSD. A non-paged \emph{reference} path, which keeps all tiles resident
and trains them in mini-batches, serves only as a yardstick.

\paragraph{What this paper reports.}
All results are from one CPU, one small corpus (TinyStories~\citep{eldan2023tinystories},
byte level), one model size (2,051,840 coefficients) and one seed (7), except a single second-seed
run of one recipe. We report:
\begin{enumerate}[leftmargin=*,itemsep=2pt]
  \item A 24-hour run with a global inverse-square-root decay of the tile learning rate, which
  plateaued above the matched-exposure bigram and spent 22\% of its wall time on paging I/O
  (Section~\ref{sec:global}).
  \item A paging change (write-behind buffer, visit-frequency eviction, larger resident set) that
  left the learning curve unchanged while reducing bytes written by about 75$\times$ between
  two 30-minute pilots, plus two ablations run on it (Section~\ref{sec:paging}).
  \item A round of 30-minute rate pilots comparing a per-tile rate floor, fixed-threshold warm
  restarts and adaptive warm restarts (Section~\ref{sec:rate}).
  \item A long run with adaptive per-tile warm restarts that passed the bigram
  (Section~\ref{sec:long}).
  \item A second pilot round ablating the restart options and trying a constant tile rate and a
  continuous plasticity law based on the Lambert $W$ function (Sections~\ref{sec:pilot2},~\ref{sec:options}).
  \item Two 6-hour runs, constant rate and Lambert plasticity, which resolve the question the restart
  run left open: up to 500k positions a constant tile rate with no restarts and no forgetting does as
  well as dynamic restarts (Section~\ref{sec:confound}); and a second seed (Section~\ref{sec:seed}).
\end{enumerate}

\paragraph{What is and is not new.}
None of the individual techniques is new: paging and offloading, mixture-of-experts style routing,
exact per-block gradients, warm restarts~\citep{loshchilov2017sgdr} and Adam~\citep{kingma2015adam}
are all established. The contribution is their combination into a working bounded-residency,
per-tile online training loop for a routed language model, and measurements of what that loop costs
and how it learns. The main empirical lesson is a simple one: in our runs this loop stalled when the tile rate
decayed on the global clock and kept learning when the rate stayed up; the restart machinery we built
to keep it up turned out to be unnecessary up to 500k positions.
We make no claim that this path is competitive with conventional training: the
non-paged reference path on the same model reached 2.216~\bpb{} at 2.09M training positions, and a
byte 5-gram fitted to about 2M bytes reaches 1.55. We also make no claim yet that the model's total
size can grow while residency stays fixed; every run here has seven tiles, and that test is future
work (Section~\ref{sec:future}).

\section{Related work}

\paragraph{Out-of-core and offloaded training.}
ZeRO-Offload and ZeRO-Infinity~\citep{ren2021zerooffload,rajbhandari2021zeroinfinity} train models
larger than accelerator memory by placing optimizer state and parameters in CPU memory or on NVMe and
streaming them through each step. They preserve the semantics of dense synchronous training, so every
parameter is read and written every step. Hexweave instead changes the training semantics so that a
step touches only a small, routed subset of parameters; the cost is that the result is no longer
dense backpropagation over the whole model.

\paragraph{Local and decoupled learning.}
Decoupled neural interfaces~\citep{jaderberg2017dni} train modules from predicted (synthetic)
gradients so that modules need not wait for a full backward pass; \citet{czarnecki2017understanding}
analyse when such predictions are adequate. Local error signals~\citep{nokland2019local}, greedy
layerwise training~\citep{belilovsky2019greedy} and the Forward-Forward
algorithm~\citep{hinton2022forward} remove or shorten the global backward pass. Earlier connected
protocols in this project (v1--v5.3, summarised in Section~\ref{sec:history}) trained tiles from a
learned, label-conditioned correction predictor in the spirit of synthetic gradients. They reached
accurate credit (prefit cosine up to 0.97) but never beat frozen tiles on held-out loss. The runs
reported here therefore use \emph{exact} gradients along each training position's own route of at
most four visits. The update is local in which parameters it touches, not in how its gradient is
computed.

\paragraph{Modular routing.}
Mixtures of experts~\citep{jacobs1991adaptive} and sparsely gated
layers~\citep{shazeer2017outrageously,fedus2021switch} route each input to a few of many experts,
so compute per example stays small while parameters grow. They are normally trained with dense
synchronous steps and all experts resident in accelerator memory. Hexweave's tiles play a similar
role, but routing is sequential (a path of visits, not a parallel top-$k$), route choices are
trained with a score-function (REINFORCE-style~\citep{williams1992reinforce}) surrogate, and expert
residency is bounded by paging.

\paragraph{Schedules, restarts and plasticity.}
Warm restarts (SGDR,~\citep{loshchilov2017sgdr}) periodically reset a decaying learning rate. Our
per-tile variant triggers a restart from accumulated relative parameter change rather than a fixed
period, and adapts the trigger from the loss of the finished cycle; in our runs a constant rate did
as well up to 500k positions (Section~\ref{sec:confound}). The baseline decay is the
inverse-square-root form used in the Transformer schedule~\citep{vaswani2017attention}. Continual
learning work studies the opposite failure, forgetting~\citep{kirkpatrick2017ewc}, and growing
capacity~\citep{rusu2016progressive,yoon2017den}. We do not address forgetting benchmarks here.

\paragraph{Prioritised replay.}
Prioritised experience replay~\citep{schaul2016per} samples stored transitions in proportion to a
priority, with importance-sampling correction; the same paper notes that purely greedy
prioritisation by error focuses on a small subset of hard samples. Our replay ablation used greedy
prioritisation without sampling or correction and showed that failure (Section~\ref{sec:paging}).

\section{The connected system}
\label{sec:system}

\subsection{Model}
\label{sec:model}

The model reads bytes (vocabulary 256) with a 512-byte window. A shared tied embedding
$E\in\mathbb{R}^{d\times 256}$, $d=128$, embeds each visible byte; each prefix row is the
RMS-normalised~\citep{zhang2019rmsnorm} sum of the byte's embedding and a sinusoidal position
vector. Seven tiles sit on a radius-1 axial hex grid ($1+3r(r+1)=7$). One \emph{anchor} is one
training position: a visible prefix and the byte that follows it. The route for an anchor starts at
the centre tile with the last prefix row as its state, and after each visit the visited tile's router
chooses a legal neighbour or halt; revisits are legal and the route stops after four visits.

A tile holds eleven matrices: single-query attention projections Q, K, V, O over the prefix
(four heads), a gated feed-forward network with gate G, input A and output DOWN at FFN width
$m=384$, a learned piecewise-linear activation curve with three knots (input clamped to the knot
range), a router, and a small \emph{pulse} that modulates the FFN output's gain and shift and adds
route biases. The pulse and router also read a retrieved \emph{history} vector: a weighted mean of
recent correction records held in a 32-record shared memory. Each visit's output is RMS-normalised.
The final state is read out through $E^\top$ plus a learned per-byte bias initialised to the
corpus' smoothed log unigram frequencies. An outer tile has 281,216 coefficients and the centre tile
282,368; the whole model has 2,051,840.

In the runs reported here, forward computation uses FP32 master weights. Each tile also stores a
ternary projection of its masters (values in $\{-1,0,+1\}$), which the project's other protocols use
for the forward pass and which we keep here only for storage and diagnostics (counting trit flips
and reversals). We make no claim about ternary kernels or ternary quality
(cf.\ \citep{ma2024era,courbariaux2015binaryconnect}).

All runs use \emph{windowed FFN slicing}: on each visit a deterministic hash of the tile identity and
coarse signatures of its input, history and prefix tail selects a circular window of 192, 288 or
336 of the 384 FFN channels; the other channels are zeroed and the DOWN projection is rescaled by
$1/\sqrt{m_{\text{active}}}$. In the 30-minute pilots that introduced it this beat the full FFN from
24,576 anchors onward (3.715 against 3.759~\bpb{} at 24,576). That is one pilot pair on one seed.

\subsection{Credit and updates}
\label{sec:credit}

For each anchor the model predicts the next byte from target-free logits, then observes the target.
An exact reverse pass along the realised route (\emph{anchor-route credit}) yields each visited
tile's parameter gradient; repeated visits to one tile are averaged. Each visited tile then takes
one Adam step on its own gradient at the tile rate $\eta_k(t)$ defined below. The shared embedding
and readout bias take one Adam step per anchor at a constant rate of 0.003. Router parameters take a
score-function update with a loss baseline~\citep{williams1992reinforce}; this is a detached surrogate
on the realised route, not an unbiased policy gradient. Weight decay is decoupled~\citep{loshchilov2019adamw}
(coefficient 0.001) and scaled by the multiplier the step actually received; after each update any
master matrix whose RMS exceeds a hard cap (1.5 or 3; Section~\ref{sec:results}) is rescaled to the cap.
Batch size is one and every byte is scored before it is learned from, so training loss is
prequential.

The update is local in the sense that matters for paging: an anchor touches at most four of seven
tiles plus the small shared state. It is not local in the stronger sense of the decoupled-learning
literature: the gradient is exact backpropagation through the route.

\subsection{The paged tile store}
\label{sec:store}

Each tile is persisted as one \emph{bundle}: FP32 masters, Adam moments (including route slots), the
packed ternary projection and per-coefficient trit-direction history. An outer bundle is 3,732,852
bytes (3.56\,MiB), dominated by training state rather than weights. The store keeps at most $R$
decoded bundles (cached plus leased) in RAM. A training step leases the bundles its route needs;
when the cache is full an idle bundle is evicted. In the original store, eviction of a dirty bundle
wrote it immediately (to a partial file, fsynced, then renamed under its content hash), and the
oldest idle bundle was evicted first.

The paging round (Section~\ref{sec:paging}) added three controls, none of which changes numerics:
\begin{itemize}[leftmargin=*,itemsep=2pt]
  \item \textbf{Write-behind buffer} of $W$ entries: an evicted dirty bundle is encoded into RAM
  instead of written. A later checkout reclaims it from RAM after a hash check; only when the buffer
  is full is the least recently evicted entry published to disk. Checkpoints drain the buffer first.
  \item \textbf{Visit-frequency eviction}: the evicted bundle is the idle one with the least
  \emph{heat}, an exponentially decayed checkout count with a half-life of about 16 checkouts.
  \item \textbf{Resident cap} $R$ raised from 4 to 6 (of 7).
\end{itemize}
A contract test checks that a bundle decoded from disk, from buffered bytes or kept cached is
bit-identical, and that losses and every tile's hash, version and step counts agree after a
checkpoint. Checkpoints (every 4,096 anchors) are complete, hash-manifested copies; the write-behind
buffer is session RAM and is never a recovery source.

The always-resident arrays (embedding, predictor kept by the protocol, memory records, one
prefix, with Adam moments) total about 1.5\,MB; the training file is streamed from disk
(Figure~\ref{fig:system}).

\begin{figure}[t]
\centering
\begin{tikzpicture}[
  font=\small,
  box/.style={draw,rounded corners=2pt,align=center,inner sep=3pt,text width=#1,
    execute at begin node=\nohyphens},
  box/.default=4.9cm,
  arr/.style={-{Stealth[length=2mm]},thick},
  lbl/.style={font=\footnotesize,fill=white,inner sep=1pt},
  panel/.style={draw=black!40,rounded corners=4pt,inner sep=5pt}
]
% Hex grid, pointy-top axial layout; IDs ordered by q then r. Dashed tile = on disk only (R = 6 example).
\foreach \q/\r/\id/\res in {-1/0/0/1,-1/1/1/0,0/-1/2/1,0/0/3/1,0/1/4/1,1/-1/5/1,1/0/6/1}{
  \pgfmathsetmacro{\x}{0.953*(\q+0.5*\r)}
  \pgfmathsetmacro{\y}{-0.825*\r}
  \ifnum\res=1
    \draw[fill=restartcol!18] (\x,\y) ++(30:0.55) -- ++(150:0.55) -- ++(210:0.55) -- ++(270:0.55) -- ++(330:0.55) -- ++(30:0.55) -- cycle;
  \else
    \draw[dashed,fill=white] (\x,\y) ++(30:0.55) -- ++(150:0.55) -- ++(210:0.55) -- ++(270:0.55) -- ++(330:0.55) -- ++(30:0.55) -- cycle;
  \fi
  \node at (\x,\y) {\id};
}
\node[font=\footnotesize,align=center] at (0,-1.85) {7 tiles, radius 1;\\ routes start at tile 3};
\node[panel,fit={(-1.5,-2.3) (1.5,1.45)},label={[font=\small\bfseries]above:Model}] (model) {};

% RAM column
\node[box,fill=restartcol!18] (res) at (6.0,1.0) {Resident set: $\le R$ decoded tile/Adam bundles (leased + idle)};
\node[box,fill=decaycol!15] (wb) at (6.0,-0.45) {Write-behind buffer: $\le W$ encoded dirty bundles};
\node[box,fill=black!5] (shared) at (6.0,-1.65) {Shared: $E$, readout bias, 32-record memory};
\node[box,fill=green!8] (tstate) at (6.0,-2.75) {Per-tile training state (outside bundles): restart clock, tile signal, plasticity};
\node[panel,fit={(res) (wb) (shared) (tstate)},label={[font=\small\bfseries]above:RAM (bounded)}] (ram) {};

% SSD column
\node[box=3.2cm] (bank) at (12.0,1.0) {Tile bank: one $\approx$3.7\,MB bundle per tile};
\node[box=3.2cm] (ckpt) at (12.0,-1.65) {Checkpoints (hash-manifested copies)};
\node[panel,fit={(bank) (ckpt)},label={[font=\small\bfseries]above:SSD}] (disk) {};

\draw[arr] (model.east |- res.west) -- node[lbl,above=1pt,align=center]{lease\\ $\le 4$ tiles} (res.west);
\draw[arr] ($(res.south)+(-1.0,0)$) -- node[lbl,left=1pt]{evict} ($(wb.north)+(-1.0,0)$);
\draw[arr] ($(wb.north)+(1.0,0)$) -- node[lbl,right=1pt]{reclaim} ($(res.south)+(1.0,0)$);
\draw[arr] (bank.west) -- node[lbl,above=1pt]{read on miss} (res.east);
\draw[arr] (wb.east) -- node[lbl,below=2pt,pos=0.45]{publish oldest} ($(bank.south)+(-0.8,0)$);
\draw[arr,dashed] (ram.east |- ckpt.west) -- node[lbl,below=1pt]{checkpoint} (ckpt.west);
\end{tikzpicture}
\caption{The connected system. Each anchor's route of at most four visits leases tile bundles into a
bounded resident set; evicted dirty bundles wait in a write-behind buffer, from which they are either
reclaimed without disk I/O or, oldest first, published to the SSD. Per-tile training state lives
outside the bundles and sets each tile's rate. The dashed tile illustrates $R=6$ of 7; the 24-hour
global-decay run used $R=4$ and no buffer.}
\label{fig:system}
\end{figure}

\subsection{Tile learning-rate schedules}
\label{sec:schedules}

All runs use base tile rate $\eta_{\text{loc}}=2\times10^{-3}$ and the shifted inverse-square-root
factor~\citep{vaswani2017attention}
\begin{equation}
  f(s) = \frac{1}{\sqrt{1+s/T}}, \qquad T = 500 \text{ anchors}.
\end{equation}

\paragraph{Global decay.} Every tile uses $\eta_k(t)=\eta_{\text{loc}}\,f(t)$, where $t$ is the global anchor
count. By $t=1.3\times10^6$ this factor is about 0.02.

\paragraph{Constant rate.} $\eta_k(t)=0.8\,\eta_{\text{loc}}=1.6\times10^{-3}$ for every tile, with
no decay and no restarts.

\paragraph{Per-tile floor.} $\eta_k(t)=\eta_{\text{loc}}\max(\varphi, f(t))$ with floor $\varphi=0.25$.

\paragraph{Per-tile warm restarts.} Each tile $k$ keeps an origin $o_k$ and uses
$\eta_k(t)=\eta_{\text{loc}}\max(\varphi, f(t-o_k))$. Each anchor produces one scalar \emph{change}
\begin{equation}
  c_t = \frac{\lVert \Delta E\rVert}{\lVert E\rVert} + \sum_{j\in\text{stepped}(t)} \frac{\lVert\Delta\theta_j\rVert}{\lVert\theta_j\rVert},
\end{equation}
the applied Adam step norm of the embedding relative to its norm, plus the same ratio for every tile
step that anchor. Every tile visited at $t$ adds $c_t$ to its \emph{drift} $D_k$ and the anchor's loss
to its current cycle's loss sum. When $D_k\ge b_k$, tile $k$ restarts: $o_k\leftarrow t$ (so its next
step runs at full rate), $D_k\leftarrow 0$, and its records are removed from shared memory so that stale
corrections are not retrieved. Adam moments are not reset. With fixed restarts the bar $b_k$ stays at the
threshold $b_0$. With \emph{dynamic} restarts and adaptation $\alpha$, from the second restart on the
finished cycle's mean anchor loss is compared with the previous cycle's: a lower loss multiplies $b_k$
by $1-\alpha$ (restart sooner), otherwise by $1+\alpha$, and $b_k$ is clamped to $[b_0/8,\,8b_0]$. We use
$b_0=2.0$, $\alpha=0.25$, $\varphi=0.25$, so the lower clamp is 0.25.

Three properties of this rule matter for interpreting the results. First, on a small fixture about 90\%
of $c_t$ came from the embedding's step, so the signal mostly measures how far a tile's \emph{inputs}
moved, and every visited tile is credited with other tiles' steps. Second, the adaptation compares
noisy single-cycle means whose anchors differ. Third, early in training it is a positive-feedback
loop that drives the bar towards its lower clamp (Section~\ref{sec:limits}).

\section{Experimental setup}
\label{sec:setup}

\paragraph{Data.} TinyStories~\citep{eldan2023tinystories} as plain text, byte level. The preparation
script terminates each story with byte \texttt{0x03} and holds out every 20th story; a SHA-256
manifest pins the files (Section~\ref{sec:repro}). Each anchor's target is the next training byte,
predicted from the previous 512 bytes. The training file is about 21\,MB and no run reported here
completed one pass.

\paragraph{Evaluation.} Every 8,192 anchors, held-out bits per byte are computed teacher-forced
with frozen weights and an isolated memory, over eight fixed public held-out spans of 128 bytes:
1,024 scored bytes per evaluation. Adjacent evaluations of one run move by about $\pm0.1$~\bpb, so
we report 12-evaluation moving means and windowed means alongside single readings, and compare runs
at matched anchors, not matched time.

\paragraph{Yardsticks.} Table~\ref{tab:yardsticks} lists reference points on this corpus. The
$n$-gram models are interpolated Witten--Bell~\citep{witten1991zero} byte models scored on the same
eight spans as the connected evaluations; the reference path used 64 spans (8,192 bytes), so compare
it by margin, not by the second decimal. The frozen-tile row is an earlier connected run whose tiles
received no updates.

\begin{table}[t]
\centering
\caption{Held-out yardsticks (\bpb, lower is better). ``Matched'' means fitted to about as many
training bytes as the connected runs had seen at comparable points.}
\label{tab:yardsticks}
\begin{tabular}{lrr}
\toprule
Model & Training bytes & Held-out \\
\midrule
Byte unigram (context-free) & -- & 4.414 \\
Connected path, tiles frozen (earlier protocol), best & 589,824 & 3.94 \\
Byte bigram, matched exposure & 625,059 & 3.31 \\
Byte trigram & 81,920 & 2.570 \\
Byte trigram & 625,059 & 2.471 \\
Byte trigram / 5-gram & 1,990,000 & 2.443 / 1.547 \\
Reference path (all tiles resident, batch 32) & 392,640 & 2.896 \\
Reference path (all tiles resident, batch 32) & 2,086,528 & 2.216 \\
\bottomrule
\end{tabular}
\end{table}

\paragraph{Hardware and process.} One desktop: 12 logical CPUs, 32\,GB RAM, NVMe SSD, Windows~11.
Every run executes under an owner script that pins the executable, places the trainer in a Job with a
512\,MiB commit cap, runs it at below-normal priority, and stops it with a checkpoint if host
available memory stays under 256\,MiB. The machine was shared with other work, and throughput varied
from 20 to 29 anchors/s for identical work.

\paragraph{Common configuration.} Seed 7 unless stated; FP32 forward; anchor-route credit; windowed FFN slicing;
base tile rate $2\times10^{-3}$ with $T=500$; shared rate 0.003; weight decay 0.001; batch one;
route length at most four. Runs differ in the resident cap $R$, write-behind size $W$, the tile-rate
schedule and the master RMS cap, as stated with each result.

\section{Results}
\label{sec:results}

\subsection{A 24-hour run with global decay}
\label{sec:global}

The first long run used global inverse-square-root decay, $R=4$, no write-behind buffer, oldest-first
eviction and RMS cap 1.5. In 24 hours it trained 1,856,133 anchors (22.6 anchors/s) with a peak
working set of 44.6\,MB. Its best single held-out evaluation was 3.271~\bpb{} at 1.74M anchors and its
last was 3.436. Windowed means (Table~\ref{tab:windows}) fell from 3.618 over the first 200k anchors
to 3.408 at 400--600k and only to 3.354 at 1.6--1.8M: a plateau that never held below the
matched-exposure bigram (Figure~\ref{fig:main}).

By the plateau the tile rate had decayed from $2\times10^{-3}$ to about $3.5\times10^{-5}$, while the
embedding and readout bias kept stepping every anchor at 0.003. All seven tiles' master RMS reached
the cap of 1.5. Visits per tile ranged from 2.12M (the centre tile, visited on every anchor) to 127k.
Of the ternary-projection direction changes, 96\% were reversals of the previous change. Our reading,
which motivated the next rounds, was that tiles were effectively frozen while their inputs kept
moving. The constant-rate run of Section~\ref{sec:confound} later supported the part of this reading
that the rate had decayed too far; this run also differs from that one in its RMS cap (1.5, which
every tile reached, against 3), a difference we did not isolate.

Paging was expensive. With four resident bundles and eager write-back the run performed 990,786 bundle
writes, wrote 3.70\,TB and read 4.6\,TB, and spent 18,229\,s, 22\% of the run, in I/O
(Table~\ref{tab:io}).

\begin{figure}[t]
\centering
\begin{tikzpicture}
\begin{axis}[
  width=0.84\linewidth, height=5.8cm, clip=false,
  xlabel={Anchors (millions)}, ylabel={Held-out \bpb},
  xmin=0, xmax=1.9, ymin=2.1, ymax=4.55, ytick={2.5,3,3.5,4,4.5},
  grid=major, grid style={black!8},
  legend style={font=\footnotesize, at={(0.99,0.99)}, anchor=north east, fill=white, fill opacity=0.9, text opacity=1},
  legend cell align=left,
  tick label style={font=\footnotesize}, label style={font=\small},
]
\addplot[only marks, mark=*, mark size=0.6pt, decaycol, opacity=0.35, forget plot]
  table[x expr=\thisrow{anchors}/1e6, y=heldout_bpb, col sep=comma]{data/global_decay_24h.csv};
\addplot[only marks, mark=*, mark size=0.6pt, restartcol, opacity=0.35, forget plot]
  table[x expr=\thisrow{anchors}/1e6, y=heldout_bpb, col sep=comma]{data/adaptive_restarts.csv};
\addplot[very thick, decaycol]
  table[x expr=\thisrow{anchors}/1e6, y=mean_bpb, col sep=comma]{data/global_decay_24h_mean.csv};
\addlegendentry{Global decay, 24 h (12-eval mean)}
\addplot[very thick, restartcol]
  table[x expr=\thisrow{anchors}/1e6, y=mean_bpb, col sep=comma]{data/adaptive_restarts_mean.csv};
\addlegendentry{Dynamic restarts, cap 3 (12-eval mean)}
\addplot[thin, black!60, dashed, domain=0:1.9, samples=2, forget plot] {4.414};
\addplot[thin, black!60, dashed, domain=0:1.9, samples=2, forget plot] {3.31};
\addplot[thin, black!60, dashed, domain=0:1.9, samples=2, forget plot] {2.57};
\addplot[thin, black!60, dotted, domain=0:1.9, samples=2, forget plot] {2.216};
\node[font=\scriptsize, anchor=west, black!70, align=left] at (axis cs:1.92,4.414) {unigram\\ 4.414};
\node[font=\scriptsize, anchor=west, black!70, align=left] at (axis cs:1.92,3.31) {bigram,\\ matched 3.31};
\node[font=\scriptsize, anchor=west, black!70, align=left] at (axis cs:1.92,2.57) {trigram\\ $\approx$2.57};
\node[font=\scriptsize, anchor=west, black!70, align=left] at (axis cs:1.92,2.216) {reference\\ path 2.216};
\end{axis}
\end{tikzpicture}
\caption{Held-out bits per byte for the two long runs. Faint points are single evaluations
(1,024 bytes each); thick lines are 12-evaluation moving means. Horizontal lines are yardsticks
from Table~\ref{tab:yardsticks}; the reference path is a different (non-paged, batched) training
path, shown only for scale. The dynamic-restart run ended at 1,691,199 anchors.}
\label{fig:main}
\end{figure}

\subsection{Paging round: write-behind and visit-frequency eviction}
\label{sec:paging}

The paging round ran 30-minute pilots with $R=6$, a three-entry write-behind buffer and
visit-frequency eviction, otherwise matching the start of the 24-hour run. The paged control
reproduced the held-out curve of the old-paging pilot to three decimals at every shared evaluation
(Table~\ref{tab:paging}), as the bit-identity contract requires, while bytes written fell from 28.9\,GB
to 0.38\,GB (about 75$\times$) and I/O time from 158\,s to 3.3\,s. Of 567 evictions, 492 were deferred
into the buffer and 480 of those were reclaimed from RAM without touching disk. The change bundles
three controls (larger $R$, buffer, eviction policy); we did not measure their separate contributions.

Two ablations ran on the new store. \emph{History off} replaces every visit's retrieved-memory input
with zeros. It made no measurable difference to held-out loss at matched anchors, but 93\% of visits
concentrated on two tiles. \emph{Greedy prioritised replay} (buffer of 1,024 past positions, two
replays per anchor, priority = most recent loss) was worse at matched anchors and ran at about half
the throughput. Replayed positions averaged 6.4 nats of loss (9.2 bits, worse than a uniform guess
over 256 bytes), consistent with the known failure of greedy
prioritisation~\citep{schaul2016per}: the buffer fixates on positions whose loss does not fall (rare
bytes, story boundaries). Replay did even out tile exposure. A replay arm with history off was also
worse.

\begin{table}[t]
\centering
\caption{Paging round, 30-minute pilots, held-out \bpb{} at matched anchors. All new-store arms use
$R=6$, $W=3$ and visit-frequency eviction. Dashes: the pilot ended before that evaluation.}
\label{tab:paging}
\begin{tabular}{lrrrr}
\toprule
Arm & 8,192 & 16,384 & 24,576 & 32,768 \\
\midrule
Old store ($R=4$, eager write, oldest-first) & 4.163 & 3.864 & 3.715 & 3.755 \\
New store, control & 4.163 & 3.864 & 3.715 & 3.755 \\
New store, history off & 3.979 & 3.822 & 3.794 & 3.777 \\
New store, greedy prioritised replay & 4.174 & 3.907 & -- & -- \\
New store, replay with history off & 4.236 & 3.941 & -- & -- \\
\bottomrule
\end{tabular}
\end{table}

\begin{table}[t]
\centering
\caption{Paging I/O. All rows are separate runs; the new-store long and 6-hour rows differ
from the 24-hour run in resident cap, buffer, eviction policy and rate schedule, so their ratios to the
24-hour run are not attributable to paging alone. Bytes per anchor are derived by division.}
\label{tab:io}
\small
\begin{tabular}{lrrrr}
\toprule
Run & Anchors & Written & I/O time & Per anchor \\
\midrule
24-hour global decay ($R=4$, eager) & 1,856,133 & 3.70\,TB & 18,229\,s (22\%) & $\approx$2.0\,MB \\
Pilot, old store ($R=4$, eager) & 30 min & 28.9\,GB & 158\,s & -- \\
Pilot, new store ($R=6$, $W=3$) & 30 min & 0.38\,GB & 3.3\,s & -- \\
Long run, new store, first 1.3M & 1,300,740 & 8.5\,GB & -- & $\approx$6.5\,kB \\
Long run, new store, whole run & 1,691,199 & 11.0\,GB & -- & $\approx$6.5\,kB \\
6-hour constant rate, new store & 509,645 & 3.45\,GB & -- & $\approx$6.8\,kB \\
6-hour Lambert, new store & 415,329 & 2.85\,GB & -- & $\approx$6.9\,kB \\
\bottomrule
\end{tabular}
\end{table}

\subsection{Rate round: floor, fixed restarts, dynamic restarts}
\label{sec:rate}

With paging cheap, a second round of 30-minute pilots on the new store changed only the tile-rate
schedule (Table~\ref{tab:rate}, Figure~\ref{fig:pilots}). A per-tile floor of 0.25 alone gave mixed
results within evaluation noise. Fixed-threshold restarts ($b_0=2.0$) were better early (3.978 against
4.163 at 8,192 anchors) but not consistently later. Dynamic restarts gave the lowest reading at all
five matched evaluations, by margins over the control (0.07--0.22~\bpb) that are only partly outside
the spread between adjacent evaluations. Dynamic restarts with the RMS cap raised from 1.5 to 3 were identical to dynamic restarts
through 40,960 anchors, so in the pilot the cap never bound.

Restart counts were very uneven. Under dynamic restarts the seven tiles restarted 103, 5, 105, 178, 6,
1 and 2 times (tile IDs 0--6; tile 3 is the centre). The three busiest tiles' bars fell from 2.0 to
0.27--0.50, near the lower clamp of 0.25, and their rate factors ended near 0.8.

\begin{table}[t]
\centering
\caption{Rate round, 30-minute pilots on the new store, held-out \bpb{} at matched anchors, and the
last evaluation of each pilot (at differing anchor counts, so not comparable across rows).}
\label{tab:rate}
\begin{tabular}{lrrrrrr}
\toprule
Arm & 8,192 & 16,384 & 24,576 & 32,768 & 40,960 & Last (anchors) \\
\midrule
Control (global decay) & 4.163 & 3.864 & 3.715 & 3.755 & 3.679$^\dagger$ & 3.661 (37,074) \\
Floor 0.25 only & 4.152 & 3.816 & 3.750 & 3.722 & -- & 3.616 (35,127) \\
Fixed restarts, $b_0=2.0$ & 3.978 & 3.730 & 3.624 & 3.734 & 3.741 & 3.558 (50,283) \\
Dynamic restarts ($\pm$25\%) & 3.941 & 3.709 & 3.595 & 3.673 & 3.607 & 3.460 (46,000) \\
Dynamic restarts, RMS cap 3 & 3.941 & 3.709 & 3.595 & 3.673 & 3.607 & 3.480 (45,886) \\
\bottomrule
\end{tabular}

\smallskip
{\footnotesize $^\dagger$From the old-store pilot, whose curve the control reproduces at every shared evaluation.}
\end{table}

\begin{figure}[t]
\centering
\begin{tikzpicture}
\begin{axis}[
  width=0.62\linewidth, height=4.5cm,
  xlabel={Anchors (thousands)}, ylabel={Held-out \bpb},
  xmin=6, xmax=52, ymin=3.4, ymax=4.3,
  grid=major, grid style={black!8},
  legend style={font=\scriptsize, at={(1.02,1)}, anchor=north west},
  legend cell align=left,
  tick label style={font=\footnotesize}, label style={font=\small},
]
\addplot[thick, black!60, mark=o, mark size=1.4pt] table[x expr=\thisrow{anchors}/1e3, y=heldout_bpb, col sep=comma]{data/pilot_control.csv};
\addlegendentry{Control}
\addplot[thick, green!50!black, mark=triangle*, mark size=1.6pt] table[x expr=\thisrow{anchors}/1e3, y=heldout_bpb, col sep=comma]{data/pilot_floor.csv};
\addlegendentry{Floor 0.25}
\addplot[thick, decaycol, mark=square*, mark size=1.2pt] table[x expr=\thisrow{anchors}/1e3, y=heldout_bpb, col sep=comma]{data/pilot_restart.csv};
\addlegendentry{Fixed restarts}
\addplot[thick, restartcol, mark=*, mark size=1.3pt] table[x expr=\thisrow{anchors}/1e3, y=heldout_bpb, col sep=comma]{data/pilot_adapt.csv};
\addlegendentry{Dynamic restarts}
\addplot[thick, purple, dashed, mark=x, mark size=1.8pt] table[x expr=\thisrow{anchors}/1e3, y=heldout_bpb, col sep=comma]{data/pilot_nohist.csv};
\addlegendentry{History off}
\addplot[thick, black!40, dotted, mark=diamond*, mark size=1.6pt] table[x expr=\thisrow{anchors}/1e3, y=heldout_bpb, col sep=comma]{data/pilot_replay.csv};
\addlegendentry{Greedy replay}
\end{axis}
\end{tikzpicture}
\caption{Thirty-minute pilots. Each point is one 1,024-byte evaluation; the last point of each pilot
falls where its 30 minutes ended. Differences below about 0.1~\bpb{} are within the spread between
adjacent evaluations of a single run.}
\label{fig:pilots}
\end{figure}

\subsection{A long run with dynamic restarts}
\label{sec:long}

The dynamic-restart recipe with RMS cap 3 then ran as a long run on the same model with $R=6$ and $W=3$.
Its windowed held-out means (Table~\ref{tab:windows}, Figure~\ref{fig:main}) fell from 3.431 over
the first 200k anchors to 2.885 at 400--600k and 2.773 at 600--800k, held near 2.78 to 1.0M, then fell
to 2.650 at 1.0--1.2M and 2.621 at 1.2--1.4M. The best single evaluation through 1.3M anchors was 2.551 at 1,269,760
anchors. From 200k anchors onward every window lies below the matched-exposure bigram (3.31); the
windows from 400k onward lie between the bigram and the trigram fitted to a similar number of bytes
(Table~\ref{tab:yardsticks}). Throughput was about 21 anchors/s, the peak working set about
50\,MB, and 8.5\,GB were written over the first 1.3M anchors.

At 1.3M anchors, restarts per tile ranged from 437 to 4,955, every bar sat at or near its lower clamp,
and rate factors were 0.56--0.95. Per-tile master RMS was 1.33--2.46, so the earlier cap of 1.5 would
have bound in this run even though it never bound in the pilot. As a derived comparison, the most
restarted tile restarted on average about once every 260 anchors, whereas the global decay factor at
1.3M anchors is about 0.02; busy tiles therefore trained at roughly 30--50 times the rate they would
have had under global decay.

The run was interrupted at 1,300,740 anchors when other programs pushed host available memory below the
owner's floor; the trainer's own peak was 54\,MB. It was resumed from its verified checkpoint. It ended at its 24-hour training budget after 1,691,199 anchors (about 20 anchors/s in the final segment). The last evaluation was 2.563 \bpb{} and the best single evaluation was 2.476 at 1,622,016 anchors; the windowed mean was 2.623 at 1.4--1.6M and 2.557 over the final 1.6--1.69M (Table~\ref{tab:windows}). Summed over its segments the run wrote about 11.0\,GB. At the end, restarts per tile ranged from 888 to 6,481, bars from 0.29 to 0.67 (lower clamp 0.25), rate factors from 0.63 to 0.94 and per-tile master RMS from 1.44 to 2.47; visits per tile ranged from 2.55M to 251k.

\begin{table}[t]
\centering
\caption{Windowed means of held-out \bpb: the mean of all evaluations whose anchor count $a$
satisfies $\text{lo}\le a<\text{hi}$, computed from the exported series in \code{paper/data/}. The
final window of each run is partial.}
\label{tab:windows}
\begin{tabular}{lrr}
\toprule
Anchors & Global decay, 24 h & Dynamic restarts, cap 3 \\
\midrule
0--200k & 3.618 & 3.431 \\
200--400k & 3.482 & 3.061 \\
400--600k & 3.408 & 2.885 \\
600--800k & 3.409 & 2.773 \\
800k--1.0M & 3.393 & 2.778 \\
1.0--1.2M & 3.375 & 2.650 \\
1.2--1.4M & 3.372 & 2.621 \\
1.4--1.6M & 3.374 & 2.623 \\
1.6--1.8M & 3.354 & 2.557 (to 1.69M) \\
1.8M--end & 3.404 & -- \\
\midrule
Best single evaluation & 3.271 (1.74M) & 2.476 (1.62M) \\
Last evaluation & 3.436 (1.86M) & 2.563 (1.69M) \\
\bottomrule
\end{tabular}
\end{table}

\paragraph{Samples.}
Generation opens with public held-out text that ends at a story boundary and fills the 512-byte window,
followed by the prompt (shown in bold), because training always places the newest byte at the last
window position. From the 1,300,740-anchor checkpoint, 200 bytes each:

\begin{quote}\small
\emph{Greedy:} \textbf{Once upon a time}, there was a little girl named Spot with her fish and she was a
big was a little box. She was said to play with her fish and said, ``Yes, she said
``Yes, I wanted to betth and said a lot her and she wa

\emph{Nucleus, $p=0.9$, $T=0.8$:} \textbf{Once upon a time}, there was a little betto named Tim whas a
smem wish. One day, Tim was very fin. She desten the to box named Amy always where st whow and do. Once
uponed to rang in the played to real with her fing.

\emph{Nucleus, $p=0.9$, $T=0.8$:} \textbf{Tom and his dog went to the} wanted to the pailed named
Specired. Amy was folled with help. One day, One day, Amy had a little beall named Spot. The became was
to the dog. He was saw the pand to the the put to help. She was a st
\end{quote}
For contrast, the 24-hour global-decay run's sample:
\begin{quote}\small
\textbf{The sun was} `` The here birt the The day haide the ware wat big boor bean it wat
was was the was foout.
\end{quote}
The restart run produces TinyStories-like words, names and dialogue structure with frequent misspelling
and little coherence; the global-decay run produces mostly function words and fragments. Each
completion took about 3.9\,s on the CPU including checkpoint load. These samples are illustrations; we
did not score them.

\subsection{Second pilot round: restart options, constant rate and Lambert plasticity}
\label{sec:pilot2}

The long run left a confound. Its bars sat at their lower clamp, so its busiest tiles restarted every
few hundred anchors and trained at 56--95\% of base rate. The intended mechanism, ``restart a tile when
its inputs have drifted, and forget its stale memory records'', and a simpler one, ``the global decay
(to about 2\% of base by 1.3M anchors) was too fast'', predict the same gain; the first pilots could
not separate them, since by 35k anchors the global factor had only fallen to $f\approx 0.12$.

The build containing the per-tile code of Section~\ref{sec:options}, run with the original
dynamic-restart configuration (RMS cap 3), reproduced the earlier pilot to three decimals at every
shared evaluation (3.941, 3.709, 3.595, 3.673, 3.607 at 8,192--40,960 anchors). That run is the
baseline of a second round in which eight 30-minute arms ($R=6$, $W=3$, RMS cap 3, seed 7) each changed
the schedule or one restart option (Table~\ref{tab:pilot2}; options in Section~\ref{sec:options}).

\begin{table}[t]
\centering
\caption{Second pilot round, 30 minutes per arm, held-out \bpb{} at matched anchors and the mean of
the five readings. ``Base'' is $\eta_{\text{loc}}=2\times10^{-3}$. Values are rounded from
\code{paper/data/pilot2\_*.csv}.}
\label{tab:pilot2}
\small
\begin{tabular}{lrrrrrr}
\toprule
Arm & 8,192 & 16,384 & 24,576 & 32,768 & 40,960 & Mean \\
\midrule
Constant rate $0.8\times$ base, no decay, no restarts & 3.900 & 3.756 & 3.570 & 3.674 & 3.459 & 3.672 \\
Dynamic restarts with \code{stagger} 64 & 3.941 & 3.717 & 3.565 & 3.692 & 3.530 & 3.689 \\
Baseline: dynamic restarts (reproduction) & 3.941 & 3.709 & 3.595 & 3.673 & 3.607 & 3.705 \\
Dynamic restarts, \code{clamp} 64 & 3.941 & 3.709 & 3.556 & 3.742 & 3.606 & 3.711 \\
Lambert plasticity, bar 0.02, floor 0.25 & 3.864 & 3.766 & 3.587 & 3.685 & --$^{\dagger}$ & 3.725$^{\dagger}$ \\
Dynamic restarts, \code{forget} off & 3.981 & 3.763 & 3.692 & 3.712 & 3.622 & 3.754 \\
Dynamic restarts, \code{replay} 4 & 3.943 & 3.790 & 3.627 & 3.734 & 3.679 & 3.755 \\
Constant rate $0.8\times$ base, restarts only clear memory & 4.020 & 3.807 & 3.675 & 3.662 & 3.632 & 3.759 \\
Dynamic restarts, \code{signal} tile, bar 0.1 & 4.048 & 3.833 & 3.718 & 3.724 & 3.675 & 3.800 \\
\bottomrule
\end{tabular}

\smallskip
{\footnotesize $^\dagger$The Lambert pilot ended at 39,291 anchors; its mean is over four readings.
The baseline's four-reading mean is 3.730.}
\end{table}

No arm beat the baseline by more than the spread between adjacent evaluations of one run
($\pm0.1$). The constant-rate and Lambert arms had the lowest means, within noise. Stagger, a
64$\times$ bar range, keeping memory records at a restart, replaying four recent positions at each
restart, and clearing memory on top of a constant rate did nothing measurable. \code{forget} off and
\code{replay} 4 trailed the baseline at all five readings but by less than the noise (mean gaps 0.049
and 0.050), and the two memory comparisons point in opposite directions (keeping records slightly
worse under restarts, clearing them slightly worse under a constant rate, mean gap 0.087), so we
conclude only that forgetting had no consistent effect at this horizon. With the 64$\times$ range the
busiest tile restarted 942 times without a change in the curve. Replayed positions had losses of
1.8--3.4 nats per tile, against 6.4 for greedy prioritised replay (Section~\ref{sec:paging}).

Only the tile-signal arm's gap is clearly outside the noise: behind at every reading by 0.05--0.12,
0.095 in the mean. Its bars collapsed from 0.1 to 0.01--0.02 (lower clamp 0.0125) and its tiles
restarted 300--2,973 times in 30 minutes, clearing memory almost continuously. We read this as a
threshold and clamp tuning failure, not as evidence against a tile-specific signal. The Lambert arm's
lifetime mean dial was 1.04--1.64 per tile, so it trained above base rate.

\subsection{Resolving the confound: 6-hour constant-rate and Lambert runs}
\label{sec:confound}

The two arms without discrete restarts then ran for 6 hours each, with the long run's configuration
otherwise ($R=6$, $W=3$, RMS cap 3, seed 7). Table~\ref{tab:sixhour} and Figure~\ref{fig:sixhour}
compare them with the two long runs at matched anchors.

\begin{figure}[t]
\begin{minipage}[b]{0.47\linewidth}
\centering
\small
\setlength{\tabcolsep}{4pt}
\begin{tabular}{lrrrr}
\toprule
Anchors & Decay & Restarts & Constant & Lambert \\
\midrule
0--100k & 3.685 & 3.546 & 3.496 & 3.496 \\
100--200k & 3.552 & 3.315 & 3.188 & 3.186 \\
200--300k & 3.514 & 3.110 & 3.062 & 3.000 \\
300--400k & 3.451 & 3.011 & 2.995 & 2.870 \\
400--500k & 3.419 & 2.926 & 2.904 & 2.871\rlap{$^3$} \\
500--600k & 3.396 & 2.842 & 2.871\rlap{$^2$} & -- \\
\bottomrule
\end{tabular}
\captionof{table}{Windowed means of held-out \bpb{} over 100k-anchor windows
($\text{lo}\le a<\text{hi}$). Windows hold 12 or 13 evaluations; a superscript gives the count where it
is smaller. Decay is the 24-hour run (RMS cap 1.5, $R=4$); the others use cap 3 and $R=6$.}
\label{tab:sixhour}
\end{minipage}\hfill
\begin{minipage}[b]{0.5\linewidth}
\centering
\begin{tikzpicture}
\begin{axis}[
  width=\linewidth, height=5.2cm,
  xlabel={Anchors (thousands)}, ylabel={Held-out \bpb{} (12-eval mean)},
  xmin=0, xmax=600, ymin=2.75, ymax=4.0, ytick={2.8,3.0,3.2,3.4,3.6,3.8},
  grid=major, grid style={black!8},
  legend style={font=\scriptsize, at={(0.99,0.99)}, anchor=north east, fill=white, fill opacity=0.9, text opacity=1},
  legend cell align=left,
  tick label style={font=\footnotesize}, label style={font=\small},
]
\addplot[very thick, decaycol, restrict x to domain=0:600]
  table[x expr=\thisrow{anchors}/1e3, y=mean_bpb, col sep=comma]{data/global_decay_24h_mean.csv};
\addlegendentry{Global decay, 24 h}
\addplot[very thick, restartcol, restrict x to domain=0:600]
  table[x expr=\thisrow{anchors}/1e3, y=mean_bpb, col sep=comma]{data/adaptive_restarts_mean.csv};
\addlegendentry{Dynamic restarts}
\addplot[very thick, green!50!black, dashed]
  table[x expr=\thisrow{anchors}/1e3, y=mean_bpb, col sep=comma]{data/const_6h_mean.csv};
\addlegendentry{Constant rate, 6 h}
\addplot[very thick, purple, dash dot]
  table[x expr=\thisrow{anchors}/1e3, y=mean_bpb, col sep=comma]{data/lambert_6h_mean.csv};
\addlegendentry{Lambert, 6 h}
\addplot[thin, black!60, dashed, domain=0:600, samples=2, forget plot] {3.31};
\node[font=\scriptsize, anchor=north east, black!70] at (axis cs:595,3.30) {bigram, matched 3.31};
\end{axis}
\end{tikzpicture}
\caption{Twelve-evaluation moving means over the first 600k anchors (seed 7), starting at the 12th
evaluation. The constant-rate run ended at 509,645 anchors, the Lambert run at 415,329.}
\label{fig:sixhour}
\end{minipage}
\end{figure}

\paragraph{Constant rate against restarts.} The constant-rate run trained 509,645 anchors at 24.9
anchors/s; last evaluation 2.818, best 2.776 at 434,176 anchors; 3.45\,GB written; per-tile master
RMS 0.96--2.27; visits per tile from 786k to 23k. In every completed window it was at or below the
restart run, by 0.050, 0.127, 0.048, 0.016 and 0.022~\bpb{} from 0--100k to 400--500k (its 500--600k
window holds two evaluations). We read this as resolving the confound for the first 500k anchors on
this seed: a constant tile rate with no restarts and no forgetting matches or slightly beats dynamic
restarts, so the warm restarts were a mechanism for keeping the rate up, not the source of the gain.
Restarts are not shown to be harmful (from 300k anchors the runs are within a few hundredths), only
unnecessary at this scale and horizon. Too-fast decay is likewise the most economical explanation of
the 24-hour run's plateau, with the RMS-cap caveat of Section~\ref{sec:global}. Nothing here tests
beyond 500k anchors: the restart run went on improving to 1.69M, and the constant-rate run was not run
that long. Each arm is one run on one seed.

\paragraph{Lambert against constant rate.} The Lambert run trained 415,329 anchors at 20.2 anchors/s;
last evaluation 2.886, best 2.765 at 352,256 anchors; 2.85\,GB written; per-tile RMS 1.57--2.81; no
discrete restarts and no memory clearing. It tied the constant-rate run over the first 200k anchors
and led it by 0.062 at 200--300k and 0.125 at 300--400k. That lead is suggestive and larger than the
noise we would expect in a 12-evaluation mean, but it is confounded exactly as the restart run was:
the dials' lifetime means were 1.07--1.45 per tile (final distances 0.025--0.049 against a bar of
0.02), so it trained at 1.07--1.45 times base against the constant-rate run's 0.8. A constant-rate
control at 1.3 times base is running and is not reported here. The Lambert run's tile exposure was
also much flatter (446k to 73k visits per tile, against 786k to 23k); we observed this but have not
established its cause. Throughput differences are within the machine's variation and not attributed.

\paragraph{Samples.} From the final checkpoints (about 0.51M and 0.42M anchors, against 1.30M for
the samples above, so not like for like; unscored; prompt in bold):
\begin{quote}\small
\emph{Constant rate:} \textbf{Once upon a time}d there, named togethey ther day. Buddy saw scared saw
and leared to sery. One day, a goonder was named flittle wery. He saw not

\emph{Lambert:} \textbf{Once upon a time}, there wooghen it it down a clearn a to the goor lown, and
Tim. Tim saw laok saw a toy saw the the mom the saw and a little cam
\end{quote}

\subsection{A second seed}
\label{sec:seed}

The dynamic-restart recipe ran for 2 hours on seed 11 (174,587 anchors). Its held-out readings at
8,192--40,960 anchors were 3.936, 3.781, 3.720, 3.763 and 3.574 (seed 7: 3.941, 3.709, 3.595, 3.673,
3.607), then, with seed 7 in brackets, 3.433 [3.359] at 57,344; 3.529 [3.391] at 98,304; 3.437
[3.393] at 139,264; 3.349 [3.342] at 155,648; and 3.314 [3.237] at 172,032; the last was 3.282 at
174,587 (\code{seed11\_2h.csv}). Seed 11 trailed by about 0.1~\bpb{} from 16k to 131k anchors (mean
gap 0.10 over those 15 evaluations, range $-0.03$ to 0.17), then closed to 0.01--0.08: the same
curve with a slower start. Routing concentrated on different tiles (211k, 127k, 100k and 53k visits on
four tiles). This is a second seed for one recipe only.

\section{Restart options and continuous plasticity}
\label{sec:options}

These mechanisms passed independent code review and were each evaluated in one 30-minute pilot
(Section~\ref{sec:pilot2}), Lambert plasticity also in one 6-hour run (Section~\ref{sec:confound}).
No restart option produced a measurable change; the tile signal, as configured, made things worse.

\paragraph{Per-tile training state.} Each tile has a small record kept apart from its weights: never
inside the paged bundle, so it can be read and updated without a checkout, and never read by
inference. It holds the restart clock (origin, drift, bar, restart count, cycle loss), a plasticity
dial multiplied into the tile rate, and a \emph{tile signal}: exponential moving averages (weight
$\max(1/n, 1/256)$ for the $n$-th sample) of the tile's input concatenated with its retrieved history
input, of its mean local gradient norm, and of the anchor loss on visited anchors. After 256 visits
(one EMA horizon) a reference snapshot of the input EMA is taken, and the tile's \emph{distance} is
$\lVert \text{ema}-\text{ref}\rVert/\lVert\text{ref}\rVert$. Unlike the shared change signal of
Section~\ref{sec:schedules}, it is specific to one tile's own inputs. The tile-signal arm and both
Lambert runs used it.

\paragraph{Restart options.} \code{forget} (default on) controls whether a restart removes the tile's
memory records. \code{signal} set to \code{tile} uses the tile distance as drift, re-snapshotting the reference at
each restart. \code{stagger} enforces a minimum number of anchors between any two restarts, restarting
the tile that has waited longest. \code{clamp} sets the adaptive bar range (default 8).
\code{replay} re-runs the restarted tile's last few live positions after a restart and gives that tile
one step per replay at its fresh rate. It is capped at 6 because each tile commit ages the tile's
memory records and gain, which expire beyond a revision lag of 8, and a visited tile already takes
two commits per anchor.

Measured (Section~\ref{sec:pilot2}): \code{forget} off, \code{stagger} 64, \code{clamp} 64 and
\code{replay} 4 did nothing measurable in 30 minutes, nor did clearing memory on a constant rate.
\code{signal} \code{tile} with bar 0.1 was worse, a tuning failure of this threshold and clamp; a
different bar, or a non-adaptive one, is untested.

\paragraph{Lambert plasticity (experimental).} A continuous alternative to discrete restarts sets each
tile's dial from its distance $d$ and bar $b$ as
\begin{equation}
  p = W_0\!\left(\frac{e\,d}{b}\right), \qquad \eta_k = \eta_{\text{loc}}\max(\varphi, p), \qquad
  \text{ref} \leftarrow \text{ref} + \min\!\left(1, \tfrac{p}{64}\right)(\text{ema}-\text{ref}),
\end{equation}
where $W_0$ is the principal branch of the Lambert $W$ function~\citep{corless1996lambert}. Since
$W_0(e)=1$, $d=b$ gives the full base rate; $d=0$ gives zero (held at the floor); larger $d$ grows
roughly logarithmically, capped at 4. The reference chases the EMA at a rate that grows with $p$, so
drift raises plasticity, which in turn pulls the reference towards the input average and lowers drift.
On a small repetitive fixture this loop settled well below the bar (dials of about 0.04--0.44 depending
on the bar).

Measured: on language modelling (bar 0.02, floor 0.25, no inverse-square-root factor) it did not
settle below the bar. Lifetime mean dials were 1.04--1.64 per tile in the pilot and 1.07--1.45 over 6
hours, so in practice the law ran tiles somewhat above base rate with per-tile variation. Its lead
over the constant-rate run at 200--400k anchors is not yet separated from that higher rate
(Section~\ref{sec:confound}).

\section{Earlier connected protocols}
\label{sec:history}

For context, and so that the negative results are on record: the connected path went through five
earlier protocols trained from \emph{predicted} credit. Eight unscored ``lanes'' per anchor ran extra
routes and updated tiles from a learned correction predictor, scaled by a Lambert-$W$ gain on the
predictor's measured quality. Each version fixed one measured failure: predictions uncorrelated with real
credit (v1); residual-stream growth that starved the real correction (v2); an alignment of the
embedding toward rank one (v3); and a predictor that, without the target byte, could only learn the
conditional expectation of the correction, which vanishes as the readout fits (v4). Conditioning the
predictor on the observed target byte's embedding (v5, the conditional form of synthetic
gradients~\citep{jaderberg2017dni}) reached prefit cosine up to 0.97, but all five v5 pilots fell
behind the frozen-tile run after 5,000--7,000 anchors. The best of these protocols stayed between the
unigram and bigram levels. Removing the lanes and using exact anchor-route credit (Section~\ref{sec:credit})
is what first broke that plateau, in a 30-minute pilot (3.573~\bpb{} at 48,398 anchors). Even predicted
credit applied only along the anchor's own route reached 3.835 in the same time, which suggests the lanes,
not the predictor's accuracy, caused the earlier plateau; that is one pilot.

\section{Discussion and limitations}
\label{sec:limits}

\paragraph{What the evidence supports.} A routed byte-level model can be trained online, one position at
a time, with only a bounded subset of its tiles decoded in RAM (peak working sets of 44.6--54\,MB for a
model whose tile bundles total about 26\,MB on disk, and whose 24-hour run read and wrote terabytes).
Paging cost depends heavily on the store policy: write-behind buffering and visit-frequency eviction
removed most writes without changing numerics. With the tile rate kept up, by per-tile warm restarts
or simply by not decaying it, the paged model passed the matched-exposure bigram and, in windowed
means, approached the reference path's 2.896 at a similar exposure (2.885 at 400--600k anchors under
restarts, 2.904 at 400--500k at constant rate, against 2.896 at 392,640), although these numbers come
from different evaluation sets and should be compared only by margin. Up to 500k anchors on seed 7,
restarts and memory forgetting were not needed for this (Section~\ref{sec:confound}).

\paragraph{The adaptive bar is a positive-feedback loop early in training.} The dynamic rule compares
consecutive cycle losses. While loss is falling overall the finished cycle nearly always improves, so
the bar shrinks at almost every restart until it reaches its lower clamp, after which the rule
effectively restarts as often as the clamp allows. The long run's bars sat at or near the clamp at
1.3M anchors, the tile-signal pilot's bars collapsed to theirs, and widening the clamp only added
restarts. An adaptive trigger would need a comparison not biased by the downward trend; we have not
implemented or tested one.

\paragraph{What it does not support.}
\begin{itemize}[leftmargin=*,itemsep=0pt,topsep=2pt]
\item \textbf{Scale independence.} Every run had seven tiles and at least four resident. With $R=6$ of
7, most of the model is resident most of the time. Whether quality holds when most tiles are on disk
is untested.
\item \textbf{Attribution beyond 500k anchors.} Whether restarts or forgetting matter after 500k
anchors, where the restart run kept improving to 1.69M, is open. The 24-hour run also differed in RMS
cap, so the size of its gap is not attributed to the decay alone.
\item \textbf{Attribution of the Lambert lead.} It coincides with a rate of 1.07--1.45 times base
against 0.8; a rate-matched control is running. The cause of its flatter tile exposure is unknown.
\item \textbf{Statistical strength.} Two seeds for the dynamic-restart recipe (over 2 hours for the
second); one seed for everything else, including the constant-rate and Lambert comparisons; one corpus;
one model size; held-out scores from 1,024 bytes that move about $\pm0.1$ between adjacent
evaluations; most comparisons are 30-minute pilots or single long runs with one run per arm.
\item \textbf{Throughput.} About 20--25 anchors/s on a shared desktop that varied 20--29 anchors/s for
identical work. The reference path ran 52--92 anchors/s with batching. The connected path is slower
per byte and has not been optimised for speed; we make no throughput claim beyond these measurements.
\item \textbf{Quality.} The best connected readings remain above a byte trigram fitted to a similar
number of bytes and far above a 5-gram. Samples are fluent-looking fragments, not coherent stories.
\item \textbf{Routing balance.} Tile exposure is skewed (2.12M to 127k visits in the 24-hour run,
786k to 23k at constant rate; 93\% of visits on two tiles without history input), and which tiles
dominate changed with the seed. The Lambert run was flatter (446k to 73k). Skew limits how much of the model each anchor
can use and will matter more with more tiles.
\end{itemize}

\section{Future work (not implemented)}
\label{sec:future}

A constant-rate control at 1.3 times base is running at the time of writing; it tests whether the
Lambert run's lead is more than its higher rate, and is not reported here. The next steps are planned
in this order, each with a pre-set gate whose failure is itself a result.
\begin{enumerate}[leftmargin=*,itemsep=0pt,topsep=2pt]
\item \textbf{Fixed larger tile count.} A radius-2 grid of 19 tiles (about 5.6M coefficients) at the
current tile size, compared with all 19 resident and with 6 resident, at matched anchors and with
4,096-byte evaluations. This is the direct test of the paper's goal: quality when most of the model
is on disk.
\item \textbf{Growth by cloning.} Split an overloaded tile into two function-preserving copies, after
Net2Net~\citep{chen2016net2net}, with weights and optimiser state copied, fresh per-tile training state
and an even router split. Splits would be triggered by a relative threshold on a per-tile signal
(measured against the other tiles, rising as reserved capacity runs out) and would use pre-allocated
padded slots so that no split waits on allocation.
\item \textbf{Batched GPU waves.} With many tiles, a wave of anchors lands mostly on different tiles
and could run as one batched operation on a GPU. An earlier GPU prototype at seven tiles ran
\emph{slower} than the CPU (17.4 against 27.9 anchors/s), so this is a hypothesis for larger tile
counts only.
\end{enumerate}

\section{Reproducibility}
\label{sec:repro}

The source, configurations and run reports are in the Hexweave repository; every number in this paper
comes from the run reports there (the plotted series are exported in \code{paper/data/}). Training
runs \code{hexweave train -{}-config <run.json>} under the owner script; samples are produced
read-only from a named checkpoint (nucleus sampling~\citep{holtzman2019curious} with the last two flags):
\begin{small}
\begin{verbatim}
cargo build --release --locked
python -B -m scripts.core_owner --config <run.json> [--resume-on-memory]
hexweave complete --config <run.json> --checkpoint <name> --prompt <text> \
    --max-new-bytes 200 [--top-p 0.9 --temperature 0.8 --seed <n>]
\end{verbatim}
\end{small}
Run names (configurations are \code{work/core-runs/<name>.json}):
\begin{footnotesize}
\begin{verbatim}
global decay, 24 h:   connected-fp32-v54-anchor-exact-window-full-s7-20261111
dynamic restarts:     connected-fp32-v56-adapt-cap3-full-s7-20261113
build reproduction:   connected-fp32-v57-repro-s7-20261116
second seed (s11):    connected-fp32-v57-adapt-cap3-s11-20261116
second pilot round:   connected-fp32-v57-<arm>-s7-20261117
constant rate, 6 h:   connected-fp32-v57-const-6h-s7-20261117
Lambert, 6 h:         connected-fp32-v57-lambert-6h-s7-20261117
corpus manifest:      2710d21277ed6e1817477ae6a61444ace674b7dc371a1ec3ea304a9e93575af8
\end{verbatim}
\end{footnotesize}
All use window 512, eight evaluation spans and evaluation every 8,192 anchors; the seed is the
\code{s7}/\code{s11} suffix. The two long runs checkpoint every 4,096 anchors. Hardware: 12 logical CPUs, 32\,GB RAM, NVMe SSD, Windows~11; training child at
below-normal priority.

\section*{Acknowledgements and use of AI tools}
Implementation, experiment execution, review and drafting of this paper were carried out with AI coding
assistants (Anthropic's Claude via Claude Code, and Google's Antigravity CLI for an earlier GPU
prototype) under the author's direction. Every reported number comes from the run reports in the
repository; the author is responsible for the content.

\bibliographystyle{plainnat}
\begin{multicols}{2}
\bibliography{refs}
\end{multicols}

\end{document}
